% ======================================================================
% 第 1 章：绪论
% 原著：R. Gupta, "Introduction to Lattice QCD" (arXiv:hep-lat/9807028)，
% 第 1 节，第 5–7 页。结尾段落延续至第 7 页（其文字摘录于 sec02.txt 顶部）。
% ======================================================================

\chapter{绪论}

本讲义关于格点QCD（LQCD），其目的在于概述既有技术层面的问题，也包括迄今在获得
具有现象学意义的数值方面所取得的进展。本讲义由三部分组成。我的职责是介绍 LQCD 并
概述 LQCD 计算的范围。在第二部分讲义中，Guido Martinelli 将讨论我们在获得结果方面
迄今所取得的进展，以及这些结果对标准模型现象学的影响。最后，Martin L\"{u}scher 将
讨论手征对称性、格点QCD 的改进形式化等热门课题，以及这些改进对下一代模拟所预期
结果的质量将产生的影响。

QCD 是强相互作用的统治性理论。它以夸克和胶子为基本表述对象，我们相信它们正是构成
强子物质的基本自由度。它在预言涉及大动量转移的现象方面非常成功。在此区域，耦合常数
很小，微扰论成为可靠的工具。另一方面，在强子世界的尺度上，$\mu \lesssim
1~\mathrm{GeV}$，耦合常数为量级一，微扰方法失效。在这个领域，格点QCD 提供了一种
非微扰工具，可以从第一性原理出发计算强子谱以及这些强子态内任意算符的矩阵元。LQCD
还可以用来处理诸如禁闭机制、手征对称性破缺、拓扑的作用，以及有限温度下 QCD 的
平衡性质等问题。遗憾的是，后面这些课题在本学校并未涉及，因此我将给出适当的参考文献
以弥补其缺失。

格点QCD 是在离散的欧几里得时空网格上表述的 QCD。由于这种离散化并未引入新的参数或
场变量，LQCD 保留了 QCD 的基本特征。格点QCD 可以有两个用途。第一，离散的时空格子
充当一种非微扰正则化方案。在有限的格距 $a$（它提供了 $\pi/a$ 处的紫外截断）下，
不存在无穷大。此外，重整化的物理量在 $a \to 0$ 时具有有限且性质良好的极限。因此，
原则上人们可以用格点正则化来完成所有标准的微扰计算，然而这些计算要复杂得多，相对于
在连续方案中进行的计算并无优势。第二，将 QCD 转写到时空格子上的首要用途是，LQCD
可以用与统计力学系统相类似的方法在计算机上进行模拟。（关于欧几里得场论与统计力学
之间联系的简要回顾见第~5 节。）这些模拟使我们能够用基本的夸克和胶子自由度来计算
强子算符的关联函数，以及任意算符在强子态之间的矩阵元。

这些模拟中唯一可调的输入参数是强耦合常数和夸克的裸质量。我们相信这些参数由某个更为
基本的底层理论所规定，然而在标准模型的框架内，它们必须由同等数量的实验量来确定。
这正是 LQCD 所做的工作。此后，如果 QCD 是强相互作用的正确理论，那么 LQCD 的所有
预言都必须与实验数据相符。

LQCD 一个非常有用的特点是人们可以调节输入参数。因此，除了检验 QCD 之外，我们还可以
对物理量关于 $\alpha_s$ 和夸克质量的依赖做出详细的预言。这些预言随后可用于约束诸如
手征微扰论、重夸克有效理论以及各种现象学模型等有效理论。

我的第一讲将致力于概述 LQCD 的范围，并说明 LQCD 的模拟是场论的分步实现。第二讲将
致力于解释如何将夸克和胶子自由度转写到格点上，以及如何构造一个在格距趋于零的极限下
给出连续 QCD 的作用量。我还将花一些时间讨论规范不变性、手征对称性、费米子加倍问题、
设计改进作用量、积分测度以及规范固定等问题。

LQCD 的数值模拟基于对欧几里得路径积分的蒙特卡洛积分，因此测量除了由格点离散化带来
的系统误差之外，还具有统计误差。为了判断格点计算的质量，重要的是理解这些误差的起源、
人们正在做些什么来量化它们，以及最终要达到给定精度的结果需要做些什么。这些问题将在
第三讲中连同对蒙特卡洛方法的基础讨论一起介绍。

第四讲将致力于 LQCD 最基本的应用——强子谱的计算，以及夸克质量和 $\alpha_s$ 的提取。
LQCD 的进展需要形式化、数值技术和计算机技术方面的改进相结合。我的总体信息是，当前
LQCD 结果——来自 $a \approx 0.05$--$0.1$~fermi 且夸克质量大致等于 $m_s$ 的模拟——
已经对标准模型现象学产生了影响。我希望这三组讲义能够成功地传达从业者的兴奋之情。

在这样一门短课程中，我只能对 LQCD 做非常简短的描述。在准备这些讲义时，我发现以下
书籍 [1--4] 和综述 [5--8] 非常有用。此外，年度格点场论会议文集 [9] 是关于当前活动
和结果状况的极好信息来源。我希望本学校的这些讲义以及这些参考文献能让任何感兴趣的
读者掌握这一主题。
